\section{Problem Formulation}
\label{sec:problem}

We formulate the decentralized routing problem in Flying Ad Hoc Networks (FANETs) as a Decentralized Partially Observable Markov Decision Process (Dec-POMDP) to formalize the separation between training-time global knowledge and execution-time local action.

\subsection{Network and Communication Model}
At any discrete time step $t$, the FANET topology is modeled as a time-varying directed graph:
\begin{equation}
    \mathcal{G}_t = (\mathcal{V}, \mathcal{E}_t),
\end{equation}
where $\mathcal{V}=\{1,\ldots,N\}$ denotes the set of UAV nodes acting as routers, and $\mathcal{E}_t \subseteq \mathcal{V} \times \mathcal{V}$ is the set of active directed wireless links. A directed link $(u, v) \in \mathcal{E}_t$ exists if node $v$ is within the communication range $R_c$ of node $u$, and the channel quality permits successful transmission at time $t$. Each UAV $u \in \mathcal{V}$ is characterized by its spatial coordinates $p_u(t) \in \mathbb{R}^3$, velocity vector $v_u(t) \in \mathbb{R}^3$, queue occupancy $q_u(t) \in [0, Q_{\max}]$, and remaining energy level $E_u(t)$. 

For a packet residing at current node $u$ destined for target node $d$, the forwarding action space $\mathcal{A}_u(t)$ is defined as the union of available one-hop neighbors and an explicit drop action:
\begin{equation}
    \mathcal{A}_u(t) = \mathcal{N}_u(t) \cup \{\dropact\},
\end{equation}
where $\mathcal{N}_u(t) = \{v \in \mathcal{V} : (u, v) \in \mathcal{E}_t\}$ represents the set of active one-hop physical neighbors of $u$, and $\dropact$ is the action corresponding to packet discarding due to buffer overflow, deadline expiration, or routing failure.

\subsection{Information Abstraction and Partial Observation}
In decentralized deployment, a local routing policy cannot access the global network graph $\mathcal{G}_t$ or the future trajectories of distant nodes. Instead, each node $u$ makes decisions based on its local observation $o_u(t)$, which maps the global state $s_t \in \mathcal{S}$ through an observation function $\mathcal{O}$:
\begin{equation}
    o_u(t) = \mathcal{O}(s_t, u) = \{x_u(t), \{x_i(t) : i \in \mathcal{N}_u(t)\}, \{x_{u,i}(t) : i \in \mathcal{N}_u(t)\}, x_d(t)\},
\end{equation}
where:
\begin{itemize}
    \item $x_u(t) = [p_u(t), v_u(t), q_u(t), E_u(t)]$ represents the self-status of node $u$.
    \item $x_i(t) = [p_i(t), v_i(t), q_i(t)]$ is the status of neighbor node $i$, periodically broadcast via lightweight single-hop beaconing.
    \item $x_{u,i}(t) = [m_i(t), \ell_i(t)]$ denotes the edge-level features of link $(u, i)$, where $m_i$ is the signal margin (RSSI) and $\ell_i$ is the predicted link lifetime.
    \item $x_d(t) = [p_d(t)]$ represents the coordinates of the destination node.
\end{itemize}

During offline training, a privileged global teacher policy $\pi_T(a \mid s_t)$ has access to the full global state $s_t$, whereas the deployable student policy $\pi_S(a \mid o_u(t))$ must select forwarding actions using only the restricted local observation $o_u(t)$.

\subsection{Routing Objective}
The primary objective of the routing policy $\pi$ is to maximize reliable, timely packet delivery to the destination while minimizing transmission delays, energy consumption, and control signaling overhead. The routing objective is formalized as:
\begin{equation}
    \max_{\pi} \; \mathbb{E}_{\pi} \left[ \sum_{t=0}^{T_{\mathrm{max}}} \gamma^t \left( R_{\mathrm{succ}} - \alpha D_t - \beta E_t - \gamma C_t - \xi H_t \right) \right],
\end{equation}
where $R_{\mathrm{succ}}$ is a positive reward received upon successful packet delivery, $D_t$ is the packet delay penalty, $E_t$ is the energy transmission cost, $C_t$ is the control footprint overhead, $H_t$ is the single-hop traversal penalty, and $\gamma \in [0, 1)$ is the discount factor.

To evaluate the proposed method against baselines, we report the following network-level and distillation-level performance metrics:
\begin{align}
\pdr &= \frac{N_{\mathrm{delivered}}}{N_{\mathrm{generated}}}, \\
\mathrm{Deadline} &= \frac{N_{\mathrm{delivered\_before\_deadline}}}{N_{\mathrm{generated}}}, \\
\bar{D} &= \frac{1}{N_{\mathrm{delivered}}} \sum_{k \in \mathcal{D}} (t_k^{\mathrm{arr}} - t_k^{\mathrm{src}}), \\
\mathrm{Delay}_{95} &= Q_{0.95} \left( \{t_k^{\mathrm{arr}} - t_k^{\mathrm{src}} : k \in \mathcal{D}\} \right), \\
\mathrm{Throughput} &= \frac{\sum_{k \in \mathcal{D}} B_k}{T_{\mathrm{sim}}}, \\
\mathrm{Overhead} &= \frac{B_{\mathrm{control}}}{N_{\mathrm{delivered}} + \epsilon}, \\
\mathrm{Energy} &= \frac{1}{N_{\mathrm{tx}}} \sum_{(u,v) \in \mathcal{E}_{\mathrm{used}}} \left( \frac{\|p_u - p_v\|_2}{R_c} \right)^2.
\end{align}
Here, $\mathcal{D}$ denotes the set of successfully delivered packets, $B_k$ is the size of packet $k$, $B_{\mathrm{control}}$ is the total control byte overhead, and $N_{\mathrm{tx}}$ is the total number of transmission attempts. The policy alignment quality is measured via the Kullback-Leibler (KL) divergence over the shared action support:
\begin{equation}
    D_{\mathrm{KL}}(\pi_T \Vert \pi_S) = \sum_{a \in \mathcal{A}_u} \pi_T(a \mid s_t) \log \frac{\pi_T(a \mid s_t)}{\pi_S(a \mid o_u(t))},
\end{equation}
and the policy return gap is quantified as:
\begin{equation}
    \Delta J = J(\pi_T) - J(\pi_S).
\end{equation}

